% ---------------------------------------------------------------------
% Section 4
% ---------------------------------------------------------------------

\begin{cnread}
本节是全文的第一次「校准」：把第 3 节的双重撕裂定理用到两个静止电荷上，
看能不能不借助任何外加的闵氏度规，直接从 \AEG{} 的几何量里读出库仑定律的
\textbf{符号规律}（同号相斥、异号相吸）。做法是：把两个电荷写成两次加法步
$q_1,q_2$，相互作用能取为切触曲率的双线性配对 $\omega(D_u,D_v)=\mu\lambda$，
于是 $\kappa=q_1q_2$ 的符号就是势能的符号；而正电子的 $+e$ 由 Paper~0
第 3.4 小节的反射算子给出，它恰好翻转了 2-链的定向（附录 D.1）。
本节 4.6 小节逐条列出哪些是已证事实、哪些是字典假设。
\end{cnread}

\subsection{Stage and dictionary}

Work on $\E_0$ and its contact model with the natural charge coordinates
$(A,M)$ of \eqref{eq:arith-contact}.  Two charges at rest are two ideal-point
histories whose multiplicative charges are equal, because both sit at the same
spatial separation from the observation point in the
Hypothesis~\ref{hyp:depth} sense,
\begin{equation}\label{eq:depth-distance}
  e^{-M_*}=\frac{\ell_0}{r}
  \qquad\Longleftrightarrow\qquad
  M_*=\log\frac{r}{\ell_0},
\end{equation}
so that $M=0$ is the ultraviolet reference (coincident limit) and $M_*\to\infty$
is infinite separation.  Charge 1 injects its additive step at the ultraviolet
end; the field it creates at the location of charge 2 is the assignment reached
after the scale step $M_*$,
\begin{equation}\label{eq:field-at-2}
  a_1(p)=q_1\,e^{-M_*}=\frac{q_1\ell_0}{r},
\end{equation}
which under \eqref{eq:depth-distance} is the Coulomb kernel up to the reference
length.

The two signs of charge are the two signs of the additive generator.

\begin{lemma}[Reflection conjugates the electron step into the positron step]
\label{lem:reflection}
Let $\mathcal R(x,y)=(-x,y)$ be the reflection of Paper~0, Subsection~3.4, which
acts on the assignment by $a\mapsto-a$, and let $\mathsf X_s$ be the
translation \eqref{eq:grid-motions}.  Then
\[
  \mathcal R\circ\mathsf X_{-e}\circ\mathcal R^{-1}=\mathsf X_{+e},
  \qquad\text{equivalently}\qquad
  a\circ\mathcal R\circ\mathsf X_{-e}\circ\mathcal R^{-1}=a+e .
\]
\end{lemma}

\begin{proof}
$\mathcal R^{-1}=\mathcal R$, and for a point $(x,y)$,
$\mathcal R(x,y)=(-x,y)$, then
$\mathsf X_{-e}(-x,y)=(-x+ey,y)$, then $\mathcal R(-x+ey,y)=(x-ey,y)
=\mathsf X_{+e}(x,y)$.
\end{proof}

\begin{warning}[The reflection is not one of the four operations]
\label{warn:reflection}
Paper~0, Subsection~3.4 is explicit that the reflection is not itself one of the
four arithmetic operations: it appears because the assignment of the model is
signed, and it is the elementary instance of the second-slot subtraction $c-z$.
Multiplication by a negative constant is the reflection composed with a positive
scaling.  Lemma~\ref{lem:reflection} must be read with that restriction: what is
conjugated is the additive motion, and the conjugation uses a map that is
outside the positive affine group.
\end{warning}

\subsection{The exchange and the interaction energy}

Consider the two histories of the exchange: charge 1 emits at the ultraviolet
end and the field is transported to charge 2, and the conjugate history in which
charge 2 emits and the field is transported back.  They are the histories
$\gamma_1$ and $\gamma_2$ of Lemma~\ref{lem:kernel-word}, with additive steps
$q_1,q_2$ in place of $p$ and multiplicative charge $M_*$ in place of $q$.

The interaction energy is then obtained in either of two equivalent ways.

\begin{enumerate}[label=(\alph*)]
  \item \emph{Pairing of field and copoint.}  The displacement of the assignment
  produced by charge 1 at the location of charge 2 is
  $\Delta a_1(r)=q_1\ell_0/r$; contracting it with the dual functional carried by
  charge 2 gives, after the $4\pi$ solid-angle normalisation,
  \begin{equation}\label{eq:pairing-energy}
    V(r)=\frac{1}{4\pi\ell_0}\,q_2\,\Delta a_1(r)
       =\frac{q_1q_2}{4\pi r} .
  \end{equation}
  \item \emph{Contact curvature density.}  Do not pair externally.  Instead let
  the additive generator strength be $\mu\propto q_1$ and the multiplicative
  coupling be $\lambda\propto q_2$.  Then the curvature scalar of Paper~0,
  Theorem~9.2 is
  \begin{equation}\label{eq:kappa}
    \omega(D_u,D_v)=\mu\lambda\;\propto\;q_1q_2=:\kappa ,
  \end{equation}
  and integrating the density over the scale interval $[0,M_*]$ selected by
  \eqref{eq:depth-distance} gives the same
  $V(r)=\kappa\,e^{-M_*}/(4\pi\ell_0)=\kappa/(4\pi r)$.
\end{enumerate}

Reading (b) is the one that carries the sign for free: by Theorem~9.2 the
curvature scalar of the horizontal distribution is literally $\mu\lambda$, and by
Remark~9.5 of Paper~0 it is an intrinsic invariant of the distribution even
though its numerical value depends on the choice of horizontal metric.  The sign
of $\kappa$ is therefore the sign of the product of the two charge steps, and
that is the whole of the sign rule.

\begin{statusbox}
\textbf{Where Hypothesis~\ref{hyp:pairing} enters.}  Theorem~9.2 supplies
$\omega(D_u,D_v)=\mu\lambda$ as an exact statement about the contact
distribution.  It does not by itself say that $\lambda$ is proportional to a
second charge, nor that the integral of the density over a scale interval is an
energy.  Both are the content of Hypothesis~\ref{hyp:pairing}, and the force law
below is conditional on it.  What is \emph{not} hypothetical is the sign
structure: once $\kappa=q_1q_2$, the two cases differ by the sign of a product
of two real numbers, and no further choice is available.
\end{statusbox}

\subsection{Like charges: the two-electron case}

Take $q_1=q_2=-e$.  Then
\begin{equation}\label{eq:kappa-ee}
  \kappa_{ee}=q_1q_2=(-e)(-e)=+e^{2}>0,
\end{equation}
which is the same sign as for two positive charges.  Geometrically this is
Lemma~\ref{lem:reflection} applied twice: the two reflections of the second
slot compose to the identity, $\mathcal R^2=\id$, so the 2-chain keeps its
orientation.  Paper~0, Appendix~D.1 makes the same statement at the level of the
oriented rectangle: the pulled-back orientation form is $pq\,\mathrm ds\wedge
\mathrm dt$, so a double sign change leaves it $+pq\,\mathrm ds\wedge\mathrm dt$.

Integrating the curvature density over the scale interval as in
\eqref{eq:pairing-energy},
\begin{equation}\label{eq:V-ee}
  V_{ee}(r)=\frac{\kappa_{ee}}{4\pi\ell_0}e^{-M_*}=+\frac{e^{2}}{4\pi r}>0 .
\end{equation}
The generalized force conjugate to the separation $r$ is
\begin{equation}\label{eq:F-ee}
  F_{ee}(r)=-\frac{\mathrm dV_{ee}}{\mathrm dr}
  =-\frac{\mathrm d}{\mathrm dr}\Bigl(+\frac{e^{2}}{4\pi r}\Bigr)
  =+\frac{e^{2}}{4\pi r^{2}}>0,
\end{equation}
i.e. directed along increasing $r$: the charges are pushed apart.

\begin{remark}[Why a positive potential means repulsion here]
\label{rem:relaxation}
The sign convention is the one fixed in Section~1.3: $V$ is the energy of the
configuration as a function of the separation, so $F=-\mathrm dV/\mathrm dr$ is
the force conjugate to $r$.  In the $\E_0$ reading of
Subsection~\ref{subsec:three-readings} below, $V_{ee}>0$ is exactly the
statement that the configuration relaxes towards $r\to\infty$, i.e. towards
$e^{-M}\to0$ on the ideal boundary.
\end{remark}

\subsection{Opposite charges: the electron--positron case}

Take $q_1=-e$ and $q_2=+e$.  By Lemma~\ref{lem:reflection} the second charge is
the first one conjugated by the reflection, so exactly one side of the 2-chain
carries a sign change.

\begin{lemma}[Orientation reversal for exactly one negative side]
\label{lem:orientation}
For the oriented singular rectangle $R_{p,q}(s,t)=(sp,tq)$ with
$0\le s,t\le1$ of Paper~0, Appendix~D.1, the pulled-back orientation form is
$pq\,\mathrm ds\wedge\mathrm dt$.  Hence the orientation of the rectangle
reverses exactly when one of $p,q$ changes sign, and the weighted integral
$\int_{R_{p,q}}e^{q-M}\mathrm dA\wedge\mathrm dM=p(e^{q}-1)$ is unchanged.
\end{lemma}

\begin{proof}
$R_{p,q}^*(\mathrm dA\wedge\mathrm dM)=pq\,\mathrm ds\wedge\mathrm dt$; the
weight is positive, so the sign of the integral is the sign of $pq$.  The value
$p(e^q-1)$ follows by direct integration as in Appendix~D.1.
\end{proof}

Consequently the curvature pairing changes sign,
\begin{equation}\label{eq:kappa-ep}
  \kappa_{e^+e^-}=q_1q_2=(-e)(+e)=-e^{2}<0,
\end{equation}
and with it the potential and the force:
\begin{equation}\label{eq:V-ep}
  V_{e^+e^-}(r)=-\frac{e^{2}}{4\pi r}<0,
  \qquad
  F_{e^+e^-}(r)=-\frac{e^{2}}{4\pi r^{2}}<0 .
\end{equation}
The force is along decreasing $r$: the charges are pulled together.

Two further \AEG{} identities describe the same case.

\begin{proposition}[The neutral pair is the closed charge loop]
\label{prop:dipole-loop}
For equal and opposite additive steps the total additive charge vanishes,
$A_{\mathrm{tot}}=(-e)+(+e)=0$, and the pair is realized by the closed
chronological loop
$\ell=(\cA_{+e},\cM_{M_*},\cA_{-e},\cM_{-M_*})$ of Paper~0, Appendix~D.2.  Its
weights satisfy
\[
  \tau(\ell,\varnothing)=e\bigl(1-e^{-M_*}\bigr)=e\Bigl(1-\frac{\ell_0}{r}\Bigr),
\]
which vanishes as $r\to\ell_0$ and grows towards $e$ as $r\to\infty$.
\end{proposition}

\begin{proof}
The charge computation is Theorem~8.6(2) with $(p,q)\mapsto(e,M_*)$; the
evaluation defect is the lemma of Section~3.5.  The limits are immediate from
\eqref{eq:depth-distance}.
\end{proof}

\begin{remark}[What the closed loop adds]
\label{rem:dipole}
The loop of Proposition~\ref{prop:dipole-loop} is the identity that Paper~0
records to separate ``zero total charge'' from ``operator neutrality''.  Read
through Hypothesis~\ref{hyp:pairing} it acquires a second reading: the dipole
has zero asymptotic charge, but because the $+e$ and $-e$ steps occur at
different depths it encloses a signed rectangle in $\ACS$, and the weighted area
of that rectangle is the defect $\tau=e(1-e^{-M_*})$.  Since
$e(1-\ell_0/r)$ increases with $r$, the configuration lowers its defect by
decreasing $r$: this is the geometric source of the attraction, stated entirely
inside \AEG{} up to the identification of $\tau$ with an energy.
\end{remark}

\subsection{Three geometric readings of the sign rule}
\label{subsec:three-readings}

\paragraph{(1) In $\E_0$: opening versus closing.}
Two electrons push the assignment in the same direction, $\mathsf X_{-e}$ each.
To keep the terminal state fixed they need compensating rescaling, and the
compensation grows with the scale span $M_*$; the pair behaves like an opening
spring whose relaxation is towards larger $r$.  Electron and positron push in
opposite directions, and at $M_*=0$ the two motions cancel identically,
$\mathsf X_{+e}\circ\mathsf M_0\circ\mathsf X_{-e}=\id$, so the configuration is
already neutral in the coincident limit; any nonzero separation inserts a scale
step $M_*$ between the two steps and produces the defect
$e(1-e^{-M_*})$ that the pair then reduces by approaching.

\paragraph{(2) In $\ACS$: divergent stacking versus a closed rectangle.}
Two like charges give $A_{\mathrm{tot}}=-2e\ne0$: no closed loop, hence no
cancellation, and the additive charge accumulates without bound as
$M_*\to\infty$ (Corollary~\ref{cor:bare-divergence}).  Opposite charges give
$A_{\mathrm{tot}}=0$ and the closed rectangle of
Proposition~\ref{prop:dipole-loop}, whose weighted area is the binding defect.

\paragraph{(3) In $\Cont$: the sign of the curvature and the handedness of the
contact volume.}
The curvature two-form is $\omega=\mu\lambda\,\mathrm du\wedge\mathrm dv$ with
$\mu\lambda\propto q_1q_2$; the contact volume form is
$\alpha\wedge\mathrm d\alpha=\mu\lambda\,\mathrm du\wedge\mathrm da\wedge
\mathrm dv$ (Paper~0, Propositions~9.1 and 9.6).  Changing the sign of exactly
one charge therefore reverses the orientation of the contact volume while
leaving the horizontal distribution unchanged in rank and integrability.  In the
reading of Theorem~\ref{thm:dual-tearing}(2) this is a reversal of the Reeb
drift: for like charges the Reeb height of the pair grows with decreasing $r$,
for unlike charges it decreases.

\subsection{Status of this section}

{\small
\renewcommand{\arraystretch}{1.22}%
\begin{longtable}{@{}p{0.55\linewidth}p{0.39\linewidth}@{}}
\toprule
statement & status \\
\midrule
\endfirsthead
\toprule
statement & status \\
\midrule
\endhead
\midrule
\multicolumn{2}{r}{\emph{continued}}\\
\endfoot
\bottomrule
\endlastfoot
$\mathcal R\circ\mathsf X_{-e}\circ\mathcal R^{-1}=\mathsf X_{+e}$
& proved (Lemma~\ref{lem:reflection}); Paper~0, Subsec.~3.4 \\
orientation of the rectangle reverses when exactly one side is negative
& proved (Lemma~\ref{lem:orientation}); Paper~0, App.~D.1 \\
$(-p)(-q)=+pq$ versus $(-p)(+q)=-pq$ for the curvature pairing
& proved once $\kappa=\mu\lambda\propto q_1q_2$; Paper~0, Props.~9.1, 9.6 \\
the neutral pair is a closed loop with defect $e(1-e^{-M_*})$
& proved; Paper~0, Thm.~8.6, App.~D.2 \\
\midrule
the additive step $\cA_{-e}$ is a physical charge $-e$, and
$\mathcal R$ is charge conjugation
& Hypothesis~\ref{hyp:charge} \\
$e^{-M_*}=\ell_0/r$ is the Coulomb kernel
& Hypothesis~\ref{hyp:depth} \\
$V=\kappa e^{-M_*}/(4\pi\ell_0)$ is an interaction energy, i.e.
$\mu\lambda$ pairing is energy
& Hypothesis~\ref{hyp:pairing} \\
$\tau=e(1-e^{-M_*})$ is a binding energy
& Hypothesis~\ref{hyp:pairing}, via Remark~\ref{rem:dipole} \\
\end{longtable}}

The three-dimensional normalisation by $4\pi$ is not derived here.  It is the
solid angle of the unit sphere, inserted because the \AEG{} model is one
dimensional in the multiplicative direction while the physical field lives in
three space dimensions; the same insertion is used in every section that quotes
a force.  A genuinely three-dimensional carrier is beyond the scope of the
present note and is recorded as an open end in Section~\ref{sec:summary}.
